
\subsection{Descriptive statistics and stylized facts}
Descriptive statistics for the SPY daily excess growth rate series confirm that the three canonical stylized facts are present in both the in-sample (2014--2024, $T=2{,}766$) and out-of-sample (2025, $T=249$) windows (Table~\ref{tab:descriptive_stats}). Excess kurtosis (kurtosis minus~3, so that a Gaussian distribution has a value of zero) reaches 7.7 (IS) and 6.9 (OoS), the Jarque-Bera normality test is rejected ($p<0.001$) in both windows, and the Ljung-Box test on $|G_t|$ is likewise rejected ($p<0.001$), confirming leptokurtosis, negligible raw-return autocorrelation, and persistent volatility clustering. Any faithful generative model must reproduce all three properties simultaneously.

% ── Table 1: SPY Descriptive Statistics ──────────────────────────────────────
\begin{table}[tp]
\centering
\caption{Descriptive statistics and stylized-facts tests for SPY daily excess growth rates.
    In-sample: 2014--2024 ($T=2{,}766$); out-of-sample: 2025 ($T=249$).
    Gaussian and Laplace columns show MLE fits to the in-sample data.
    JB = Jarque-Bera normality test.
    LB = Ljung-Box autocorrelation test at lag 20.
    $^\ast$ denotes rejection at $\alpha=0.05$.}
\label{tab:descriptive_stats}
\small
\begin{tabular}{@{}lcccc@{}}
\toprule
\textbf{Statistic} & \textbf{IS Observed} & \textbf{OoS Observed} & \textbf{Gaussian} & \textbf{Laplace} \\ \midrule
Mean (annualized, \%)       & 6.31   & 11.60  & 6.31   & 14.19  \\
Std Dev (annualized, \%)    & 214.50 & 220.62 & 214.46 & 205.93 \\
Skewness                    & $-0.753$ & $-1.093$ & 0    & 0      \\
Excess Kurtosis             & 7.715  & 6.867  & 0      & 3      \\
JB normality test           & reject$^\ast$ ($p<0.001$) & reject$^\ast$ ($p<0.001$) & n/a   & n/a   \\
LB test on $G_t$ (lag 20)   & reject$^\ast$ ($p<0.001$) & reject$^\ast$ ($p=0.019$) & n/a   & n/a   \\
LB test on $|G_t|$ (lag 20) & reject$^\ast$ ($p<0.001$) & reject$^\ast$ ($p<0.001$) & n/a   & n/a   \\ \bottomrule
\end{tabular}
\end{table}


\subsection{Jump hyperparameter optimization and state resolution}
A multi-objective grid search over $\epsilon\in[10^{-4},\,2.5\times10^{-2}]$ and $\lambda\in[10,\,160]$, with 200 simulated paths per grid point, identified a well-defined interior minimum at $\epsilon^{*}=10^{-4}$ and $\lambda^{*}=100$, indicating that the jump hyperparameters are identifiable and that the optimization has a clear, unique minimum (Figure~\ref{fig:parameter_sweep}, Online Appendix~S2). At these optimal values, the ACF of $|G_t|$ closely tracks the observed SPY ACF across lags 1--252, confirming that the jump mechanism reproduces the empirical ARCH effect.

The primary analysis uses $N=100$ states, but a sensitivity study over $N\in\{30,60,90,100,150,200\}$ shows that the framework is robust to the choice of state resolution across a wide range (Table~\ref{tab:sensitivity_results}, Online Appendix~S3). KS and AD pass rates remain stable for both HMM-NJ ($\geq$98.9\% KS, $\geq$97.9\% AD) and HMM-WJ ($\geq$96.0\% KS, $\geq$88.6\% AD), while ACF-MAE for HMM-WJ improves slightly from 0.057 at $N=30$ to 0.049 at $N=200$. State-occupancy diagnostics confirm that every state at $N\leq200$ contains at least 7 observations, ensuring adequate statistical support per state under the Laplace equal-probability partition. At $N=350$, however, certain states are never visited in the historical record, leaving zero transition probabilities that prevent the model from reaching all states, establishing a practical upper bound on state resolution.


\subsection{In-sample distributional and temporal quality}
We generated 1,000 synthetic paths of 2,766 trading days from each of six generators and evaluated them against the observed SPY series using two-sample KS and AD tests at $\alpha=0.05$, along with kurtosis matching, ACF fidelity, novelty, diversity, and quantile coverage (Table~\ref{tab:model_comparison}).

\paragraph{Baseline anchors.}
We included three non-structured baselines to bracket the quality landscape. Bootstrap resampling of the training data serves as a non-parametric upper bound: we observed 100.0\% KS and 99.7\% AD pass rates with 100\% quantile coverage, but this generator simply recycles observed values and cannot produce genuinely new distributional structure (0 estimated parameters). The Gaussian i.i.d.\ baseline (MLE fit) achieved 0.0\% KS and 0.0\% AD pass rates, which confirms that a symmetric thin-tailed model is categorically inadequate for heavy-tailed equity returns. The Laplace i.i.d.\ baseline fared better at 44.0\% (SE~1.6) KS and 43.3\% (SE~1.6) AD, which shows that matching the tail decay rate alone recovers roughly half the distributional quality; the remaining gap reflects the regime-dependent variation that a single Laplace distribution cannot capture.

\paragraph{Distributional fidelity.}
We evaluated KS/AD pass rates and kurtosis matching for the three structured models (Table~\ref{tab:model_comparison}). GARCH(1,1) achieved IS KS and AD pass rates of only 5.5\% (SE~0.7) and 1.9\% (SE~0.4), which indicates that the GARCH marginal distribution is systematically inconsistent with the observed data despite its well-known volatility clustering properties. HMM-NJ achieved 99.4\% (SE~0.2) and 98.6\% (SE~0.4), and HMM-WJ achieved 97.0\% (SE~0.5) and 90.5\% (SE~0.9), which confirms that the Laplace-mixture emission structure adequately captures the observed heavy-tailed distribution. Simulated excess kurtosis tells a complementary story: the observed value is 7.7, GARCH overshoots at 8.2 (SE~0.37) with high variance across paths, while HMM-NJ and HMM-WJ produce 5.7 (SE~0.02) and 5.5 (SE~0.03) respectively. The HMM variants underestimate kurtosis because the equal-probability Laplace partition smooths the most extreme tail quantiles, but they do so with negligible path-to-path variation, which suggests a stable generative mechanism. Quantile coverage reinforces these findings: Bootstrap, HMM-NJ, and HMM-WJ all achieve 100\% coverage of the 1st--99th empirical percentiles, while Gaussian covers only 13.1\%, Laplace 37.4\%, and GARCH 29.3\%, which indicates systematic tail misspecification in all three non-HMM parametric generators (Figure~\ref{fig:model_comparison}, panel~c).

\paragraph{Temporal fidelity.}
We measured temporal quality using the mean absolute error of $\text{ACF}(|G_t|)$ over lags 1--252. Bootstrap, Gaussian, and Laplace all produce ACF-MAE of 0.060, which is expected because i.i.d.\ draws destroy all temporal dependence by construction. GARCH achieved the lowest ACF-MAE (0.031, SE~0.001) because its variance equation is explicitly designed to match autocorrelation structure; however, this advantage comes at the cost of distributional fidelity, a tradeoff that the quality metrics expose directly. HMM-NJ scored 0.059 (SE${}<$0.001), essentially indistinguishable from the i.i.d.\ floor, which confirms that the standard transition matrix alone is structurally incapable of generating persistent volatility clustering (Figure~\ref{fig:model_comparison}, panel~b). HMM-WJ scored 0.052 (SE${}<$0.001), the only non-GARCH model to reduce ACF-MAE below the i.i.d.\ baseline. The improvement arises because approximately 24\% of HMM-WJ paths contain at least one Poisson jump event that sustains $\text{ACF}(|G_t|)$ decay across all lags, while the remaining 76\% behave identically to HMM-NJ. The jump-duration mechanism is therefore the structural feature responsible for reproducing the ARCH effect.

\paragraph{Novelty, diversity, and parsimony.}
We computed novelty (mean correlation distance to the observed series) and diversity (mean pairwise correlation distance among synthetic paths) to verify that the generators produce genuinely new, non-redundant scenarios. All six models scored between 0.984 and 0.985 on both metrics, which confirms that no generator memorizes the training data and that synthetic paths are mutually independent. This uniformity is expected: all generators produce stochastic paths that are uncorrelated with any specific historical realization. Finally, HMM-WJ achieves the best joint distributional and temporal quality among all parametric models while requiring only four estimated parameters ($N$, $\epsilon$, $\lambda$, and the transition matrix $\mathbf{P}$, where $N$ is set by design), compared with three for GARCH and two each for the Gaussian and Laplace baselines.

% ── Table 2: Model Comparison ───────────────────────────────────────────────
\begin{table}[tp]
\centering
\caption{In-sample and out-of-sample model comparison for SPY ($N=100$, 1{,}000 simulated paths,
    significance level $\alpha=0.05$).
    Bootstrap: i.i.d.\ resample of training data.
    Gaussian/Laplace: i.i.d.\ draws from MLE-fitted distributions.
    GARCH: GARCH(1,1) estimated by quasi-maximum likelihood.
    HMM-NJ: standard Hidden Markov Model without jumps.
    HMM-WJ: hybrid model with Poisson jump-duration extension.
    ACF-MAE: mean absolute error of the ACF of $|G_t|$ over lags 1--252.
    Novelty: mean correlation distance $1-|\rho|$ between each synthetic path and the observed series.
    Diversity: mean pairwise correlation distance among synthetic paths.
    Coverage: fraction of 99 empirical quantiles (1st--99th) falling within the
    [5th, 95th] percentile envelope of the corresponding synthetic quantiles.
    Standard errors in parentheses: binomial SE for pass rates,
    $\text{std}/\!\sqrt{n}$ for kurtosis and novelty,
    bootstrap SE ($B=500$) for ACF-MAE, diversity, and coverage.
    Arrows indicate preferred direction: $\uparrow$~higher is better,
    $\downarrow$~lower is better, $\approx$~closer to observed is better.}
\label{tab:model_comparison}
\resizebox{\columnwidth}{!}{%
\footnotesize
\begin{tabular}{@{}lcccccc@{}}
\toprule
\textbf{Metric} & \textbf{Bootstrap} & \textbf{Gaussian} & \textbf{Laplace} & \textbf{GARCH(1,1)} & \textbf{HMM-NJ} & \textbf{HMM-WJ} \\ \midrule
\multicolumn{7}{l}{\textit{In-sample: 2{,}766 trading days (2014--2024)}} \\[2pt]
KS pass rate (\%) $\uparrow$            & 100.0 ({$<$}0.1) &   0.0 ({$<$}0.1) &  44.0 (1.6) &   5.5 (0.7) &  99.4 (0.2) &  97.0 (0.5) \\
AD pass rate (\%) $\uparrow$            &  99.7 (0.2)      &   0.0 ({$<$}0.1) &  43.3 (1.6) &   1.9 (0.4) &  98.6 (0.4) &  90.5 (0.9) \\
Excess kurtosis (observed)   & \multicolumn{6}{c}{7.715} \\
Excess kurtosis (simulated) $\approx$  &   7.6 (0.08) & $-$0.0 ({$<$}0.01) &   3.0 (0.02) &   8.2 (0.37) &   5.7 (0.02) &   5.5 (0.03) \\
ACF-MAE $\downarrow$                      & 0.060 ({$<$}0.001) & 0.060 ({$<$}0.001) & 0.060 ({$<$}0.001) & 0.031 (0.001) & 0.059 ({$<$}0.001) & 0.052 ({$<$}0.001) \\
Novelty $\uparrow$                      & 0.985 ({$<$}0.001) & 0.984 ({$<$}0.001) & 0.985 ({$<$}0.001) & 0.985 ({$<$}0.001) & 0.984 ({$<$}0.001) & 0.984 ({$<$}0.001) \\
Diversity $\uparrow$                    & 0.985 ({$<$}0.001) & 0.985 ({$<$}0.001) & 0.985 ({$<$}0.001) & 0.985 ({$<$}0.001) & 0.984 ({$<$}0.001) & 0.984 ({$<$}0.001) \\
Coverage (\%) $\uparrow$                & 100.0 ({$<$}0.1)   &  13.1 (0.1)       &  37.4 (1.5) &  29.3 (1.5) & 100.0 ({$<$}0.1)   & 100.0 ({$<$}0.1) \\[4pt]
\multicolumn{7}{l}{\textit{Out-of-sample: 249 trading days (2025)}} \\[2pt]
KS pass rate (\%) $\uparrow$            &  97.4 (0.5) &  62.2 (1.5) &  88.0 (1.0) &  80.3 (1.3) &  96.5 (0.6) &  95.0 (0.7) \\
AD pass rate (\%) $\uparrow$            &  98.5 (0.4) &  46.3 (1.6) &  92.9 (0.8) &  72.0 (1.4) &  96.8 (0.6) &  95.9 (0.6) \\
Excess kurtosis (observed)   & \multicolumn{6}{c}{6.867} \\
Excess kurtosis (simulated) $\approx$  &   6.0 (0.18) & $-$0.0 ({$<$}0.01) &   2.7 (0.05) &   1.6 (0.07) &   5.0 (0.08) &   4.9 (0.08) \\
ACF-MAE $\downarrow$                      & 0.043 ({$<$}0.001) & 0.043 ({$<$}0.001) & 0.043 ({$<$}0.001) & 0.026 ({$<$}0.001) & 0.041 ({$<$}0.001) & 0.040 ({$<$}0.001) \\
Novelty $\uparrow$                      & 0.949 (0.001) & 0.950 (0.001) & 0.952 (0.001) & 0.951 (0.001) & 0.948 (0.001) & 0.949 (0.001) \\
Diversity $\uparrow$                    & 0.950 ({$<$}0.001) & 0.949 ({$<$}0.001) & 0.949 ({$<$}0.001) & 0.950 ({$<$}0.001) & 0.948 ({$<$}0.001) & 0.948 ({$<$}0.001) \\
Coverage (\%) $\uparrow$                & 100.0 ({$<$}0.1) &  36.4 (1.9) &  74.7 (1.1) &  96.0 (1.3) & 100.0 ({$<$}0.1) & 100.0 ({$<$}0.1) \\[4pt]
\midrule
Parameters estimated         &     0 &     2 &     2 &     3 &     2 &     4 \\ \bottomrule
\end{tabular}}
\end{table}


% ── Figure 4: Head-to-Head Model Comparison ──────────────────────────────────
\begin{figure}[tp]
    \centering
    \includegraphics[width=\textwidth]{Fig4-Model-Comparison.pdf}
    \caption{Head-to-head in-sample model comparison for SPY ($N=100$, 1{,}000 simulated paths).
    Panel~(a): marginal density of excess growth rates with IS KS pass rates annotated.
    GARCH(1,1) fails the two-sample KS test on 95\% of paths (pass rate 5.5\%), indicating a
    systematic shape mismatch with the empirical distribution.  Both HMM variants pass at
    rates $\geq\!97\%$, confirming that the Laplace-mixture emission adequately captures the observed
    heavy tails.  HMM-WJ generalises robustly out-of-sample (OoS KS pass rate 95\%), while
    GARCH recovers partially (OoS 80\%).
    Panel~(b): autocorrelation function of absolute excess growth rates,
    $\mathrm{ACF}(|G_t|)$, at lags 1--252 (shaded bands: 10th--90th percentile across paths).
    HMM-NJ is structurally incapable of producing persistent volatility clustering:
    without a jump mechanism the latent states are i.i.d.\ conditional on the Markov chain,
    so $\mathrm{ACF}(|G_t|)\approx 0$ for all lags beyond one (dotted curve).
    HMM-WJ generates a \emph{mixture} of two path families under the same model:
    approximately 76\% of IS paths contain no Poisson jump and therefore behave
    identically to HMM-NJ (dashed curve, near zero); the remaining $\sim$24\% of paths
    contain at least one jump event and sustain substantial $\mathrm{ACF}(|G_t|)$ decay
    across all lags (solid navy curve $\pm$ band).
    The jump frequency, and hence the fraction of the ensemble exhibiting slow ACF decay, is
    controlled by the hyperparameter $\lambda$, making the volatility-clustering strength
    directly tunable.
    Panel~(c): tail Q-Q plot at the 0.1st--99.9th percentile region; HMM-WJ provides the
    closest mean quantile match in the extreme tails.}
    \label{fig:model_comparison}
\end{figure}

\subsection{Out-of-sample evaluation}
We evaluated all six generators on a held-out window of 249 trading days (January~2 through December~31, 2025) using the same metrics as the in-sample analysis (Table~\ref{tab:model_comparison}). Because the OoS window is roughly one-ninth as long as the training set, all two-sample tests have lower statistical power; pass rates therefore rise mechanically across every model, and the meaningful comparison is relative ranking rather than absolute level.

\paragraph{Baselines and distributional fidelity.}
We observed that Bootstrap resampling achieves OoS KS and AD pass rates of 97.4\% (SE~0.5) and 98.5\% (SE~0.4), remaining near the ceiling. The Gaussian baseline rises from 0\% IS to 62.2\% (SE~1.5) KS and 46.3\% (SE~1.6) AD, and Laplace from 44\% IS to 88.0\% (SE~1.0) KS and 92.9\% (SE~0.8) AD; these improvements reflect the reduced test power, not genuine distributional fit, as confirmed by their kurtosis values (Gaussian 0.0, Laplace 2.7 against an observed 6.9). GARCH recovers to 80.3\% (SE~1.3) KS and 72.0\% (SE~1.4) AD, but its simulated kurtosis collapses to 1.6, which indicates that GARCH fails to preserve tail behaviour outside the training window. HMM-NJ achieves 96.5\% (SE~0.6) and 96.8\% (SE~0.6), and HMM-WJ achieves 95.0\% (SE~0.7) and 95.9\% (SE~0.6), with simulated kurtosis near 4.9--5.0, confirming that both HMM variants generalise robustly.

\paragraph{Coverage, temporal fidelity, novelty, and diversity.}
Quantile coverage separates the models: Bootstrap, HMM-NJ, and HMM-WJ maintain 100\%, GARCH recovers to 96.0\% (from 29.3\% IS), and Laplace reaches 74.7\%, while Gaussian remains low at 36.4\%. OoS ACF-MAE values are lower across the board because the shorter window reduces autocorrelation estimation precision: Bootstrap, Gaussian, and Laplace all score 0.043, GARCH achieves 0.026, HMM-NJ 0.041, and HMM-WJ 0.040. The gap between HMM-NJ and HMM-WJ narrows from 0.007 IS to 0.001 OoS, which reflects reduced ACF estimation power rather than convergence in temporal quality. Novelty and diversity remain uniformly high across all models (0.948--0.952), confirming that no generator memorises the training data even in a shorter evaluation window.

Nonetheless, the out-of-sample period exposes the stationarity assumption embedded in the IS-calibrated jump parameters. The distribution of KS $p$-values shifts leftward from IS to OoS, reflecting the expected distributional degradation when stationary parameters are applied to a period with elevated macro uncertainty (Figure~\ref{fig:statistical_validation}, panels~a--b). The observed ACF of $|G_t|$ in the 2025 window exceeds the HMM-WJ simulation band at medium-to-long lags, indicating that $(\epsilon^{*},\lambda^{*})$ underestimate the frequency and persistence of volatility episodes in this period and that the jump parameters are regime-dependent (Figure~\ref{fig:statistical_validation}, panel~d). Despite this ACF gap, the observed marginal density remains within the simulation envelope, confirming adequate distributional coverage even under regime shift (Figure~\ref{fig:statistical_validation}, panel~c).

% ── Figure 6: Statistical Validation ─────────────────────────────────────────
\begin{figure}[tp]
    \centering
    \includegraphics[width=\textwidth]{Fig6-Statistical-Validation.pdf}
    \caption{Statistical validation of HMM-WJ ($N=100$, 1{,}000 simulated paths, $\alpha=0.05$).
    Panels~(a) and~(b) show the distribution of two-sample KS $p$-values in-sample and
    out-of-sample, respectively; a well-calibrated generative model produces $p$-values that are
    approximately uniform above the significance threshold.  The leftward shift from~(a) to~(b)
    reflects the expected distributional degradation when stationary parameters are applied to a
    regime with elevated macro uncertainty.
    Panel~(c) shows the out-of-sample marginal density fan chart; the observed density (dashed red)
    falls within the 10th--90th percentile simulation envelope, confirming adequate distributional
    coverage despite the regime shift.
    Panel~(d) shows the ACF of $|G_t|$ in the out-of-sample window.  The observed ACF (dashed red)
    exceeds the simulation band at medium-to-long lags, indicating that the 2025 test period
    exhibited stronger volatility clustering than the IS-calibrated jump parameters
    $(\epsilon^*,\lambda^*)$ predict.  This gap provides direct evidence that the jump
    hyperparameters are regime-dependent and motivates the time-varying extensions discussed in
    Section~\ref{sec:discussion}.}
    \label{fig:statistical_validation}
\end{figure}


\subsection{Multi-asset extension}
Propagating HMM-WJ SPY paths through a Single-Index Model for 424 S\&P~500 constituents, with $\hat{G}_{i,t}=\hat{\alpha}_i+\hat{\beta}_i G_{{\rm SPY},t}+\hat{\eta}_{i,t}$ estimated by OLS on the training window and idiosyncratic shocks drawn by resampling empirical residuals, yields a median IS KS pass rate of 66.7\% and a mean of 58.4\% across the universe (Figure~\ref{fig:multi_asset}, panel~b; Table~\ref{tab:multi_asset}). These rates are lower than the 97.0\% achieved for SPY alone and reflect the well-known limitation that a single-factor linear decomposition with symmetric residuals cannot fully capture asset-specific tail behaviour, skewness, or sector-level dynamics. The SIM $R^2_i$ varies widely across GICS sectors, from near-zero for defensive and low-correlation names to above 0.8 for broad-market ETFs (Figure~\ref{fig:multi_asset}, panel~a), and the degradation in KS pass rate concentrates among assets with low $R^2_i$. High-$\beta$ assets closely tracking SPY retain pass rates above 80\%, suggesting that the HMM-WJ generative engine is not the bottleneck; the limiting factor is the expressiveness of the factor model (Figure~\ref{fig:multi_asset}, panel~c).

Out-of-sample KS pass rates across 417 surviving assets reach a mean of 82.1\% and a median of 91.8\%, confirming that the SIM extension generalises to the test window (Table~\ref{tab:multi_asset}). The correlation-distance distribution between simulated and historical paths yields a mean distance of $1-\rho\approx1.0$, confirming that the framework produces temporally independent realisations rather than tracking specific historical events, consistent with its role as a scenario generation tool rather than a forecasting instrument (Figure~\ref{fig:statistical_validation}).

% ── Table 4: Multi-Asset SIM Extension Results ──────────────────────────────
\begin{table}[tp]
\centering
\caption{Cross-sectional summary statistics for the Single-Index Model (SIM) extension
    across S\&P~500 constituents.  Individual asset paths are generated as
    $\hat{G}_{i,t}=\hat{\alpha}_i+\hat{\beta}_i G_{{\rm SPY},t}+\hat{\eta}_{i,t}$
    where $G_{{\rm SPY},t}$ is drawn from HMM-WJ ($N=100$) simulated paths
    and $\hat{\eta}_{i,t}$ is resampled from empirical residuals.
    KS pass rates computed over 1{,}000 paths at $\alpha=0.05$.}
\label{tab:multi_asset}
\small
\begin{tabular}{@{}lcccc@{}}
\toprule
\textbf{Statistic} & \textbf{Mean} & \textbf{Median} & \textbf{5th pct} & \textbf{95th pct} \\ \midrule
\multicolumn{5}{l}{\textit{In-sample (424 assets, 2{,}766 trading days, 2014--2024)}} \\[2pt]
KS pass rate (\%)                    & 58.4   & 66.7   &  2.5   & 96.0   \\
SIM $R^2$                            &  0.298 &  0.295 &  0.107 &  0.507 \\
$\hat{\beta}$ (factor loading)       &  1.066 &  1.086 &  0.480 &  1.657 \\
$\hat{\alpha}$ (intercept)           & $-0.0328$ & $-0.0196$ & $-0.2037$ & 0.0964 \\[4pt]
\multicolumn{5}{l}{\textit{Out-of-sample (417 assets, 249 trading days, 2025)}} \\[2pt]
KS pass rate (\%)                    & 82.1   & 91.8   & 23.8   & 99.0   \\
SIM $R^2$                            &  0.142 &  0.142 & $-0.169$ &  0.479 \\
$\hat{\beta}$ (factor loading)       &  1.064 &  1.084 &  0.482 &  1.660 \\
$\hat{\alpha}$ (intercept)           & $-0.0320$ & $-0.0186$ & $-0.2035$ & 0.0965 \\ \bottomrule
\end{tabular}
\end{table}